% ============================================================================
% CAPÍTULO 5 — DISCUSSÃO
% ============================================================================
\chapter{Discussão}
\label{cap:discussao}

\section{Alinhamento com a literatura de referência}

\begin{table}[htb]
    \caption{Comparação dos achados com os artigos-base (Experimento A)}
    \label{tab:alinhamento-literatura}
    \centering
    \resizebox{\textwidth}{!}{%
    \begin{tabular}{lcccc}
        \toprule
        Traço & De Bortoli (2019) & Jiang et al. (2024) & Este estudo (Exp. A) & Análise de robustez \\
        \midrule
        Abertura à Experiência & Sig. positivo & Sig. positivo & Sig. positivo & Estável \\
        Extroversão            & Não sig.       & Otimismo (indireto) & Sig. positivo & Estável \\
        Consciência            & Não sig.       & Sig. negativo & Sig. positivo & Instável \\
        Neuroticismo           & Não sig.       & Sig. negativo & Sig. positivo & Instável \\
        Amabilidade             & Não sig.       & Papel mínimo & Não sig. & Sem efeito consistente \\
        \bottomrule
    \end{tabular}%
    }
    \legend{Fonte: elaborado pela autora a partir de \cite{debortoli2019}, \cite{jiangpengyan2024} e das Tabelas \ref{tab:ols} e \ref{tab:robustez-sintese} (Capítulo \ref{cap:resultados}).}
\end{table}

O efeito positivo e significativo de Abertura à Experiência sobre o
escore do IPA (Tabelas \ref{tab:spearman-ipa} e \ref{tab:ols}) é o
resultado mais robusto deste trabalho, por ser o único traço que replica
de forma consistente nos dois artigos-base. Em \citeonline{debortoli2019},
Abertura à Experiência foi o único traço estatisticamente significativo
para prever o perfil de risco real de estudantes em uma simulação de
investimento; em \citeonline{jiangpengyan2024}, o mesmo traço está entre
os dois com maior poder explicativo sobre o comportamento de investidores
reais, com efeito robusto sobre a alocação em ações nas três amostras do
estudo (Estados Unidos, Austrália e Alemanha). Os autores deste segundo
artigo decompõem o mecanismo por trás desse efeito: ao contrário do
Neuroticismo, cujo efeito passa pela leitura pessimista do cenário
econômico, o efeito de Abertura à Experiência não depende de crenças
sobre o mercado: decorre de uma disposição mais direta a aceitar risco,
independentemente de como o investidor avalia as perspectivas econômicas
(seção \ref{cap:referencial}).

Esse padrão, o mesmo traço na mesma direção em três contextos
metodologicamente distintos, é uma evidência relevante de robustez,
embora não uma prova. As três amostras diferem em praticamente tudo que
poderia gerar um resultado espúrio: sujeitos (estudantes brasileiros;
investidores reais americanos, majoritariamente homens brancos com mais
de 60 anos e patrimônio elevado; agentes de linguagem artificial),
variável dependente (perfil obtido por simulação de investimento real;
alocação declarada de carteira; escore de um questionário de perfil de
investidor respondido pelo próprio agente) e método estatístico
(regressão logística ordenada; OLS; correlação de Spearman e OLS). Um
efeito que se sustenta apesar dessa variação tem menor probabilidade de
ser um artefato de amostra, de instrumento ou de escolha estatística
específica de um dos três estudos. Ao mesmo tempo, essa convergência não
deve ser lida como evidência de que o agente de IA ``possui'' uma
disposição psicológica real ao risco equivalente à humana; ela mostra,
de forma mais modesta, que o pipeline de indução e mensuração deste
trabalho é capaz de reproduzir, ao menos para este traço, uma relação já
bem estabelecida na literatura, o que funciona também como uma validação
indireta do desenho experimental adotado, sem prejuízo das limitações de
validade externa desse tipo de simulação discutidas adiante neste
capítulo.

\textbf{Extroversão} não foi confirmada por \citeonline{debortoli2019}
(traço não significativo naquele estudo) e está apenas parcialmente
alinhada com \citeonline{jiangpengyan2024}, que encontram Extroversão
associada a expectativas mais otimistas sobre o mercado, mas sem efeito
significativo sobre a alocação real em ações; ou seja, o mecanismo que
Jiang et al. relatam para Extroversão é indireto e não se traduz em
comportamento, enquanto neste trabalho o efeito sobre o escore do IPA é
direto e significativo (coeficiente 1,19, p $<$ 0,001, Tabela
\ref{tab:ols}). Duas observações, feitas antes da análise de robustez,
relativizavam esse achado. Primeiro, na seleção do modelo, o efeito não
apareceu no \texttt{llama3.2:3b}, em que a manipulação de Extroversão foi
fraca ($\rho$=0,185; Tabela \ref{tab:comparacao-modelos}), o que era
compatível tanto com um achado real mascarado por manipulação ruim quanto
com um efeito restrito ao \texttt{llama3.1:8b}. Segundo, o
Experimento B (traço isolado, seção \ref{sec:exp-a-b}) não confirma este
achado com significância estatística ($\rho$=0,242, p=0,070, N=57, Tabela
\ref{tab:ipa-exp-b}), embora a magnitude seja parecida à do Experimento A (0,340),
o que sugere que a perda de significância pode refletir apenas o N bem
menor por traço no desenho isolado, não a ausência do efeito (discutido
em detalhe na seção \ref{sec:discussao-exp-b}). A análise de robustez e
generalização resolve essa dúvida a favor do efeito: com N=300 em cada
rodada, Extroversão foi positiva e significativa ($p<0,001$) nas quatro
temperaturas, com outro sorteio de perfis e nos três modelos abertos
adicionais (Tabela \ref{tab:robustez-sintese}), além de aparecer na mesma
direção nos dois modelos Nemotron. Isso responde às duas observações: o
efeito não é restrito ao \texttt{llama3.1:8b}, e sua ausência no
Experimento B é atribuível ao N menor. Por isso, Extroversão é tratada neste
trabalho como um efeito \textbf{robusto}, ao lado de Abertura à
Experiência, ainda que sem o respaldo dos dois artigos-base que Abertura
tem.

\textbf{Consciência e Neuroticismo}, diferentemente do que ocorria no
estudo piloto (BFI-10, onde nenhum dos dois era significativo), emergiram
como \textbf{significativos e positivos} na análise principal com BFI-44
(coeficientes 0,75 e 0,41 respectivamente, Tabela \ref{tab:ols}). Isso
muda a natureza da divergência em relação à literatura de referência: já
não é uma ausência de efeito (que seria consistente com
\citeonline{debortoli2019}, onde os dois também não são significativos),
e sim um efeito \textbf{na direção oposta} à relatada por
\citeonline{jiangpengyan2024}, que encontram Consciência e Neuroticismo
significativos e \textbf{negativos} sobre a alocação em ações, de forma
robusta nas três amostras do artigo (Estados Unidos, Austrália e
Alemanha; ver Capítulo \ref{cap:referencial}, seção "Síntese e hipóteses
de referência"). Ou seja, este trabalho diverge de \emph{ambos} os
artigos-base para esses dois traços: de De Bortoli et al. por encontrar
efeito onde eles não encontraram, e de Jiang et al. por encontrar o sinal
oposto ao deles. Não há, neste trabalho, dados que permitam decidir entre
duas explicações concorrentes para esse padrão: (i) o mecanismo humano
que Jiang et al. descrevem para Neuroticismo (leitura mais pessimista
do cenário econômico, que reduz alocação em ações) pode simplesmente
não se manifestar da mesma forma em um agente de LLM, cujas "crenças"
sobre o mercado não são geradas pelos mesmos processos cognitivos
descritos na literatura de finanças comportamentais; ou (ii) a descrição
de personalidade usada no \textit{prompt} para Consciência e Neuroticismo
(perfis\_personalidade.json) pode estar sendo interpretada pelo modelo de
um modo que favorece respostas mais decididas/arrojadas no IPA: por
exemplo, se o modelo associa "alta Consciência" (organização, disciplina,
orientação a objetivos) a uma postura mais assertiva de investimento, em
vez de à cautela esperada pela literatura humana. O Experimento B (seção
\ref{sec:discussao-exp-b}) mostra que, isolados, nenhum dos dois traços
mantém significância estatística, e Neuroticismo inverte de sinal
novamente ($\rho$=$-$0,116, Tabela \ref{tab:ipa-exp-b}). A análise de
robustez e generalização confirma essa fragilidade com amostras de mesmo
tamanho que a do Experimento A: com outro sorteio de perfis, os dois
coeficientes ficam próximos de zero e sem significância (0,048 e
$-$0,092, Tabela \ref{tab:robustez-seed}); entre temperaturas, oscilam
bastante de tamanho; e entre modelos, Consciência chega a inverter o
sinal (negativa no \texttt{qwen2.5:7b}) e cada um dos dois perde a
significância em pelo menos um modelo (Tabela \ref{tab:robustez-sintese}).
Por isso, Consciência e Neuroticismo são considerados neste trabalho
\textbf{resultados instáveis}: os efeitos aparecem no Experimento A, mas
dependem do sorteio, da temperatura e do modelo, e não devem ser
interpretados como uma relação geral entre esses traços e o perfil de
investidor simulado. Essa conclusão também relativiza a divergência com
\citeonline{jiangpengyan2024} discutida acima: sem um efeito estável, o
sinal oposto ao da literatura humana não pode ser tratado como um
achado substantivo sobre os agentes.

\textbf{Amabilidade} é o único traço sem nenhuma divergência a discutir:
não apresentou efeito significativo em nenhum dos três datasets deste
trabalho (piloto, Experimento A ou Experimento B, isolado ou não), o que
é consistente com os dois artigos-base: Amabilidade também não é
significativa em \citeonline{debortoli2019} e tem "papel mínimo" em
\citeonline{jiangpengyan2024}. Dos cinco traços investigados, é o único
em que os quatro estudos (os dois artigos-base e os dois desenhos deste
trabalho) concordam integralmente. Na análise de robustez, Amabilidade
aparece significativa em algumas rodadas, mas com sinais opostos
(positiva nas temperaturas 0,0 e 0,3, negativa no \texttt{qwen2.5:7b} e
no \texttt{mistral:7b}, Tabela \ref{tab:robustez-sintese}), o que reforça
a leitura de que não há um efeito consistente desse traço sobre o escore
do IPA.

\section{Indução simultânea vs. isolada: o Experimento B}
\label{sec:discussao-exp-b}

O Experimento B foi motivado por uma pergunta metodológica levantada
durante a orientação deste trabalho: será que induzir os 5 traços do Big
Five simultaneamente no mesmo \textit{prompt}, o desenho do Experimento
A, herdado da forma como a personalidade é medida em sujeitos humanos nos
dois artigos-base, que sempre possuem os 5 traços presentes em algum
nível, dilui ou confunde o sinal de cada traço individual, em
comparação a manipular apenas 1 traço por vez? O Experimento B testa essa
possibilidade diretamente, mantendo modelo e instrumento fixos e variando
apenas o desenho de indução (seção \ref{sec:exp-a-b}).

A resposta, com os dados deste trabalho, é dividida em duas partes. Na
\textbf{checagem de manipulação}, ou seja, na capacidade do modelo de
incorporar corretamente o nível de personalidade pedido, os dois
desenhos produzem resultados praticamente idênticos (Tabela
\ref{tab:manipulacao-exp-b}, Figura \ref{fig:exp-ab-manipulacao}): a correlação entre nível induzido e
autorrelato varia entre 0,81 e 0,93 tanto com os 5 traços variando juntos
quanto com apenas 1 variando isoladamente, sem diferença sistemática a
favor de nenhum dos dois desenhos. Isso é evidência direta \textbf{contra}
a hipótese de que a indução simultânea prejudica a capacidade do modelo
de incorporar cada traço: se houvesse esse tipo de interferência, seria
razoável esperar uma checagem de manipulação mais fraca no Experimento A
(5 traços competindo por "atenção" no mesmo \textit{prompt}) do que no B
(1 traço isolado), o que não se observa.

Já na \textbf{ligação com o escore do IPA} (Tabela \ref{tab:ipa-exp-b},
Figura \ref{fig:exp-ab-ipa}), o resultado é mais heterogêneo por traço. Abertura à Experiência não só
se mantém significativa isolada como fica mais forte ($\rho$=0,792 contra
0,530 no Experimento A); o achado mais robusto deste trabalho continua
robusto sob o desenho alternativo. Os outros quatro traços, no entanto,
perdem significância estatística quando isolados, incluindo os três que
emergiram como novidade no Experimento A (Extroversão, Consciência e
Neuroticismo, seção anterior). A explicação mais provável para essa
perda não é interferência entre traços, já descartada pela
checagem de manipulação idêntica nos dois desenhos, e sim
\textbf{poder estatístico}: o Experimento B tem, para cada traço
individual, entre 49 e 69 agentes (só os agentes em que aquele traço foi
sorteado como dominante), contra 300 no Experimento A. Para Extroversão,
por exemplo, a magnitude do efeito é semelhante nos dois desenhos (0,340
no A, 0,242 no B); a diferença é se esse coeficiente atinge o limiar de
significância com N=300 ou com N=57, não necessariamente se o efeito
existe. Neuroticismo é a exceção mais chamativa, invertendo de sinal
entre os dois desenhos (positivo no A, negativo no B, Tabela
\ref{tab:ipa-exp-b}); com N=69, essa inversão é compatível com
instabilidade de amostra pequena, mas também reforça que esse é o traço
com evidência menos consistente entre todos os testados neste trabalho.

Em síntese, o Experimento B responde à pergunta original de forma
específica: \textbf{induzir os 5 traços simultaneamente não parece
prejudicar a qualidade da indução em si}, mas o desenho isolado, do jeito
que foi implementado aqui (N pequeno por traço), \textbf{não tem poder
estatístico suficiente para confirmar ou refutar com segurança} os
achados secundários do Experimento A. Essa confirmação veio da análise
de robustez e generalização, discutida na próxima seção, que manteve
N=300 em cada rodada: Extroversão se confirmou, e Consciência e
Neuroticismo não.

\section{Robustez e generalização dos resultados}
\label{sec:discussao-robustez}

A análise de robustez e generalização (seção \ref{sec:resultados-robustez})
separa com clareza dois grupos de resultados. O primeiro é formado por
Abertura à Experiência e Extroversão, que elevam o escore do IPA em todas
as oito rodadas com N=300 (quatro temperaturas, dois sorteios de perfis e
quatro modelos abertos), sempre com $p<0,001$, e também nos dois modelos
Nemotron. Como essas rodadas variam, um de cada vez, o ruído da geração,
os perfis sorteados e o modelo, é pouco provável que esses dois efeitos
sejam produto de uma configuração específica do Experimento A. O segundo
grupo é formado pelos três traços de efeito menor (Consciência,
Neuroticismo e Amabilidade), cujos coeficientes mudam de significância e
de sinal entre rodadas. A temperatura, em particular, afeta mais o ruído
do que a direção dos efeitos: reduzi-la aumenta a fidelidade da
manipulação e o R², sem alterar quais traços têm efeito estável.

A generalização entre modelos, porém, vale para a \textit{direção} dos
dois efeitos estáveis, não para o \textit{nível} de risco declarado pelos
agentes. Com os mesmos perfis de personalidade, o perfil médio de
investidor vai de majoritariamente Moderado (\texttt{llama3.1:8b}) a
quase todo Arrojado (\texttt{mistral:7b}), e o mesmo agente recebe o
mesmo perfil do \texttt{llama3.1:8b} em apenas 20,3\% a 62,0\% dos casos,
conforme o modelo (Tabela \ref{tab:robustez-modelos-pareada}). O tamanho
dos efeitos também varia: o coeficiente de Abertura à Experiência vai de
0,632 (\texttt{mistral:7b}) a 4,643 (\texttt{qwen2.5:7b}). Na prática,
isso significa que uma simulação desse tipo pode indicar
\textit{quais} traços tornam um agente mais ou menos tolerante ao risco,
mas o perfil de investidor absoluto que ela atribui a uma personalidade
é, em grande parte, uma característica do modelo usado. Esse ponto é
relevante para qualquer uso de agentes de IA como substitutos de
respondentes humanos em avaliações de \textit{suitability}.

\section{Comparação com estudos de personalidade induzida em LLMs}
\label{sec:discussao-llm}

Entre os estudos revisados na seção \ref{sec:ref-llm-risco}, o de
\citeonline{borman2024} tem interseção direta com este trabalho. O
desenho de indução é essencialmente o mesmo do Experimento A: os cinco
traços do Big Five induzidos simultaneamente por instrução, cada um em
três níveis (Baixo, Médio e Alto), com a personalidade verificada por um
teste do próprio modelo dos cinco fatores antes da análise. A diferença
está no modo de compor a amostra (os autores usaram uma persona para cada
uma das 243 combinações possíveis; este trabalho sorteou os níveis de 300
agentes) e no modelo (ChatGPT 4.0, proprietário, contra, na análise
principal e nas rodadas com N=300, modelos abertos de 7 a 12 bilhões de
parâmetros). Os dois estudos também usam regressão
OLS para estimar a relação entre os traços e o comportamento. Essa
proximidade torna a comparação mais direta do que com os artigos-base,
que trabalham com sujeitos humanos.

A comparação mais informativa é com o questionário comportamental de
\citeonline{borman2024}, já que o IPA também é um questionário. Naquele
estudo, o questionário reproduziu a direção esperada pela literatura
humana para o apetite ao risco em Abertura à Experiência, Extroversão e
Neuroticismo, mas não em Consciência e Amabilidade. Neste trabalho,
Abertura e Extroversão aparecem na direção esperada e de forma estável,
convergindo com os autores; Consciência é instável e Amabilidade não
apresenta efeito consistente, o que também converge com a dificuldade
relatada por eles para esses dois traços no questionário. A divergência está em
Neuroticismo, que \citeonline{borman2024} reproduziram na direção
esperada e que, neste trabalho, é instável e, quando significativo,
positivo. Os autores relatam ainda que as personas representaram melhor
o comportamento humano em uma tarefa simulada de investimento do que no
questionário, inclusive acertando a direção dos cinco traços no apetite
ao risco. Esse resultado sugere que parte da instabilidade observada
neste trabalho nos traços de efeito menor pode estar ligada ao formato
de questionário, e motiva uma das recomendações de trabalho futuro
(Capítulo \ref{cap:conclusao}).

O estudo de \citeonline{hartley2025} converge com este trabalho em dois
pontos, apesar das diferenças de desenho (um traço induzido por vez,
modelos proprietários e risco medido por escolhas entre loterias). O
primeiro é o papel central de Abertura à Experiência: no GPT-4o, foi o
único traço com correlação significativa com os parâmetros de
sensibilidade ao risco, assim como é, neste trabalho, o traço de maior e
mais estável efeito sobre o IPA. O segundo é a dependência do modelo: os
autores observaram que a relação entre personalidade e risco não se
manteve nas variantes menores dos modelos e chegou a se inverter em
alguns deles, o que tem paralelo na variação entre modelos encontrada
aqui, tanto na seleção do modelo (seção \ref{sec:selecao-modelo}) quanto
na análise de generalização. Diferentemente de Hartley et al., porém,
este trabalho encontrou Extroversão com efeito estável em todos os
modelos testados, o que pode refletir a diferença de instrumento
(questionário de \textit{suitability} contra loterias) ou de forma de
indução (cinco traços simultâneos contra um traço por vez).

\section{O papel da checagem de manipulação}

Além da comparação entre desenhos de indução (seção anterior), a
comparação entre os 4 modelos de linguagem testados (Tabela
\ref{tab:comparacao-modelos}, seção \ref{sec:selecao-modelo}) é, em si,
um resultado deste trabalho: mostra que a qualidade da indução de
personalidade não é uniforme entre modelos, e que essa qualidade, não
apenas o tamanho do modelo isoladamente, é pré-requisito para
interpretar qualquer efeito sobre o IPA. O caso mais claro é Abertura à
Experiência: não significativa no \texttt{llama3.2:1b} ($\rho\approx$0,09
n.s.) nem no \texttt{llama3.2:3b} ($\rho$=0,077 n.s.) na checagem de
manipulação, e por consequência também não significativa na relação com
o IPA nesses dois modelos, mas fortemente incorporada no
\texttt{llama3.1:8b} ($\rho$=0,739 no piloto BFI-10; 0,913 no Experimento
A com BFI-44), onde o efeito sobre o IPA finalmente emerge com
significância. A interpretação mais direta é que os modelos menores não
"desconfirmam" a relação entre Abertura e perfil de investidor: eles
simplesmente falham em incorporar o traço de forma detectável, tornando
impossível testar a relação com esses dados. Sem a checagem de
manipulação, um resultado não significativo em um modelo pequeno poderia
ser lido (incorretamente) como evidência contra a hipótese de pesquisa,
quando na verdade reflete uma falha na etapa de indução, anterior à
pergunta de pesquisa em si.

\section{Vieses de resposta dos agentes}

Os quatro modelos testados na seleção do modelo apresentam vieses de resposta ao IPA que
independem da personalidade induzida. O mais consistente é o
\textbf{viés de centralidade}: os modelos da família Llama concentram
respostas no perfil Moderado, em proporção que diminui à medida que o
modelo cresce: aproximadamente 96\% no \texttt{llama3.2:1b}, 90\% no
\texttt{llama3.2:3b} e 78\% no \texttt{llama3.1:8b} (estudo piloto,
BFI-10). No Experimento A (BFI-44), o mesmo modelo (\texttt{llama3.1:8b})
concentrou 82,3\% dos agentes no perfil Moderado (Tabela
\ref{tab:distribuicao-ipa}), proporção semelhante à do piloto, o que
sugere que o viés de centralidade é uma característica do modelo, não do
instrumento de checagem de manipulação usado. O \texttt{gemma3:4b}
apresenta um viés distinto, de \textbf{conservadorismo}: no piloto
reduzido com esse modelo (N=30), 60\% dos agentes foram classificados
como Conservador e nenhum como Arrojado, mesmo com personalidades
induzidas variadas. A análise de generalização mostrou que cada modelo
tem um viés próprio, e não apenas os modelos menores: com os mesmos
perfis de personalidade, o \texttt{mistral:7b} concentrou 91,5\% dos
agentes no perfil Arrojado, o \texttt{gemma3:12b} teve a maior proporção
de Conservadores (24,0\%), e os dois modelos Nemotron responderam de
forma polarizada entre Conservador e Arrojado (Tabelas
\ref{tab:robustez-modelos} e \ref{tab:robustez-nemotron}).

Esses vieses são relevantes por dois motivos. Primeiro, como limitação de
validade: um modelo com viés de centralidade forte comprime a variância
do escore do IPA, o que reduz o poder estatístico disponível para
detectar qualquer efeito real da personalidade sobre esse escore,
prejudicando sistematicamente os modelos menores duas vezes (manipulação
mais fraca e menor variância na variável dependente). Segundo, como
achado metodológico por si só: a atenuação do viés de centralidade com o
aumento do modelo (96\%→90\%→78\%→82,3\%, aproximadamente) sugere que
não é só a capacidade de incorporar personalidade que melhora com modelos
maiores, mas também a disposição do modelo a dar respostas mais
extremas/decididas em geral, um comportamento de LLMs como sujeitos
simulados que merece atenção em trabalhos futuros que usem esse tipo de
questionário de múltipla escolha com agentes de IA.

\section{Limitações do estudo}

\begin{itemize}
    \item \textbf{O instrumento de autorrelato da análise principal não
    corresponde exatamente ao de nenhum dos dois artigos-base}. O BFI-10,
    usado no estudo piloto e na comparação entre modelos, é o mesmo
    instrumento (e a mesma tradução) usados por
    \citeonline{debortoli2019}; o BFI-44, usado na análise principal
    (Experimentos A e B), em uma versão em português que não é idêntica
    à adaptação validada para o Brasil por \citeonline{andrade2008}, não
    corresponde exatamente nem ao BFI-10
    de De Bortoli et al. nem ao SAPA-20 de \citeonline{jiangpengyan2024}.
    Essa é uma limitação real de comparabilidade introduzida pela
    migração de instrumento (seção \ref{sec:exp-a-b}), mitigada, mas
    não eliminada, pelo fato de os itens do BFI-44 medirem os mesmos 5
    traços do mesmo modelo teórico (Big Five), com escala de resposta
    idêntica à do BFI-10;
    \item \textbf{Vieses de resposta do modelo}, sobretudo a
    concentração de respostas no perfil Moderado, discutidos na seção
    anterior, comprimem a variância da variável dependente e reduzem o
    poder estatístico disponível para detectar efeitos reais da
    personalidade;
    \item \textbf{Os cortes de classificação do perfil IPA
    (Conservador/\allowbreak Moderado/\allowbreak Arrojado) são uma operacionalização própria
    deste trabalho}, não publicada no artigo original de referência para
    o questionário (seção \ref{sec:ipa-pontuacao}); outra
    escolha de corte poderia produzir uma distribuição de perfis
    diferente, embora não afete as correlações de Spearman nem os
    coeficientes da OLS, calculados sobre o escore contínuo;
    \item \textbf{Validade externa}: os sujeitos deste trabalho são
    agentes de linguagem artificial, não investidores reais: os
    achados dizem respeito a como um modelo de linguagem específico
    responde quando instruído a simular uma personalidade, não
    necessariamente a como pessoas reais com essas personalidades se
    comportariam;
    \item \textbf{Abrangência dos modelos testados}: a análise de
    generalização cobriu quatro modelos abertos de 7 a 12 bilhões de
    parâmetros com N=300 e dois modelos maiores apenas em piloto (N=30),
    acessados por uma cota gratuita de API. Modelos proprietários de
    grande porte (como os das famílias GPT, Claude e Gemini) não foram
    testados, e um modelo previsto para o piloto (Gemma 4 31B) não pôde
    ser incluído por indisponibilidade do provedor gratuito; os
    resultados podem não se estender a esses modelos;
    \item \textbf{Questionário em vez de tarefa comportamental}: o perfil
    de investidor foi medido por um questionário de autorrelato (IPA), e
    não pelo comportamento do agente em uma tarefa de investimento. Em
    \citeonline{borman2024}, personas de LLM representaram melhor o
    comportamento humano em uma tarefa simulada do que em um
    questionário, o que sugere que o formato do instrumento pode ter
    contribuído para a instabilidade dos efeitos menores;
    \item \textbf{N pequeno por traço no Experimento B} (49 a 69 agentes,
    contra 300 no Experimento A), como discutido na seção
    \ref{sec:discussao-exp-b}, isso limita a capacidade desse experimento
    de confirmar ou refutar com segurança os efeitos secundários do
    Experimento A, papel que coube à análise de robustez e
    generalização.
\end{itemize}
